\documentclass[11pt,a4paper]{article}
\usepackage[utf8]{inputenc}
\usepackage[margin=1in]{geometry}
\usepackage{amsmath,amssymb}
\usepackage{graphicx}
\usepackage{hyperref}
\usepackage{booktabs}
\usepackage{listings}
\usepackage{color}
\usepackage{tikz}
\usetikzlibrary{shapes.geometric, arrows}

\hypersetup{
    colorlinks=true,
    linkcolor=blue,
    filecolor=magenta,      
    urlcolor=cyan,
    pdftitle={MarketPulse Platform Documentation},
    pdfpagemode=FullScreen,
}

\definecolor{codegreen}{rgb}{0,0.6,0}
\definecolor{codegray}{rgb}{0.5,0.5,0.5}
\definecolor{codepurple}{rgb}{0.58,0,0.82}
\definecolor{backcolour}{rgb}{0.95,0.95,0.92}

\lstdefinestyle{mystyle}{
    backgroundcolor=\color{backcolour},   
    commentstyle=\color{codegreen},
    keywordstyle=\color{magenta},
    numberstyle=\tiny\color{codegray},
    stringstyle=\color{codepurple},
    basicstyle=\ttfamily\footnotesize,
    breakatwhitespace=false,         
    breaklines=true,                 
    captionpos=b,                    
    keepspaces=true,                 
    numbers=left,                    
    numbersep=5pt,                  
    showspaces=false,                
    showstringspaces=false,
    showtabs=false,                  
    tabsize=2
}
\lstset{style=mystyle}

\title{\textbf{MarketPulse: A Real-Time Quantitative Trading \& MLOps Data Platform}}
\author{\textbf{Intern A} \\ Department of Platform Engineering}
\date{July 16, 2026}

\begin{document}

\maketitle
\tableofcontents
\newpage

\section{Project Overview}
MarketPulse is a global markets quantitative data platform engineered to ingest, process, transform, and serve stock and cryptocurrency financial data at scale. The platform operates on a medallion architecture model, utilizing a hybrid structure that combines real-time streaming ingestion with scheduled batch transformations and machine learning prediction pipelines.

\subsection{Objectives}
\begin{itemize}
    \item Ingest real-time asset pricing data from public WebSockets/REST endpoints.
    \item Orchestrate a Bronze $\rightarrow$ Silver $\rightarrow$ Gold Medallion Pipeline in AWS S3 and Databricks.
    \item Compute technical analysis indicators (RSI, MACD, Bollinger Bands) and macroeconomic metrics.
    \item Train an LSTM Regressor model to project asset return vectors and log experiments to an MLflow Model Registry.
    \item Deploy an end-to-end serverless data serving layer utilizing Amazon Redshift Serverless and Amazon QuickSight.
\end{itemize}

\section{System Architecture}
The system divides compute and storage layers to minimize cost while maximizing analytical query performance. 

\subsection{Medallion Storage Layers (S3)}
Data is organized into three distinct structural layers in AWS S3:
\begin{enumerate}
    \item \textbf{Bronze (Raw):} Captures raw JSON and CSV tick metrics directly from brokers without any schema enforcement.
    \item \textbf{Silver (Cleaned):} Cleansed tables containing deduplicated records, corrected null values, and validated transaction dates.
    \item \textbf{Gold (Analytics-Ready):} Joined tables containing daily asset values, macro variables, calculated technical indicators, and feature vectors.
\end{enumerate}

\subsection{Data Flow Schema}
Ingested events move sequentially through:
\[
\text{Alpaca API} \xrightarrow{\text{Kinesis Stream}} \text{Bronze (S3)} \xrightarrow{\text{PySpark Cleaning}} \text{Silver (S3)} \xrightarrow{\text{dbt compilation}} \text{Gold (S3)} \xrightarrow{\text{Redshift Loader}} \text{BI / Serving}
\]

\section{Key Technical Decisions \& Workarounds}
Building this production-grade system required overcoming several platform incompatibilities:

\subsection{dbt Core \& Python 3.14 Meta Compatibility}
During deployment, compiling dbt models under pre-release Python 3.14 caused a fatal crash in Pydantic V1 metaclass fields. We resolved this by patching the virtual environment validator module (\texttt{pydantic/v1/validators.py}) to dynamically evaluate PEP 649 deferred annotations and handle \texttt{FieldInfo} parameters gracefully.

\subsection{AWS Redshift Serverless Network Provisioning}
Redshift Serverless requires subnet configurations across at least three unique Availability Zones (AZs) in a VPC to create a workgroup. Since the sandbox account lacked a default VPC, we wrote a dynamic provisioning script (\texttt{setup_redshift.py}) that automatically discovers the custom VPC, allocates two extra subnets in \texttt{us-east-1b} and \texttt{us-east-1c}, links them to the public route table, and creates the serverless workgroup.

\subsection{Delta-RS S3 File Bypassing}
To access S3 Delta tables from Databricks Serverless clusters without JVM class loader conflicts, we utilized the Rust-backed \texttt{deltalake} engine (\texttt{DeltaTable}) to handle reading and writing Delta format transactions directly from S3.

\section{Orchestration \& Data Quality}
Workflow orchestration is managed locally via Apache Airflow running in Docker containers, configured with root access to map the host's Docker socket (\texttt{/var/run/docker.sock}) to execute operator containers.

Data quality is checked at the gold layer using a customized execution script resembling Great Expectations syntax. It checks:
\begin{itemize}
    \item Non-null requirements on critical fields (\texttt{symbol}, \texttt{timestamp}).
    \item Value constraints (e.g., closing price $> 0.01$).
    \item Unique constraints to prevent duplicate database loads.
\end{itemize}

\section{BI Dashboard \& ML Models}
The ML pipeline trains an LSTM neural network on the Gold features to predict 5-day return vectors. The model registry and experiment runs are logged in \texttt{mlruns.db}. 

Redshift Serverless serves as the analytical backend, queryable directly by Amazon QuickSight and Streamlit. The QuickSight dashboard features:
\begin{enumerate}
    \item A bar chart of top buy/sell recommendations by predicted returns.
    \item A line chart displaying the ML portfolio equity curve vs. the benchmark.
\end{enumerate}

\end{document}
